\documentclass[11pt]{article}
\usepackage[margin=1in]{geometry}
\usepackage[T1]{fontenc}
\usepackage[utf8]{inputenc}
\usepackage{lmodern}
\usepackage{microtype}
\usepackage{array}
\usepackage{booktabs}
\usepackage{enumitem}
\usepackage{hyperref}
\usepackage{xcolor}
\hypersetup{
  colorlinks=true,
  linkcolor=black,
  urlcolor=blue!50!black,
  pdftitle={Coding Scheme Dossier}
}
\setlist[itemize]{leftmargin=1.6em}
\setlist[description]{style=nextline,leftmargin=3.5cm,labelsep=0.6em}
\newcommand{\field}[1]{\nolinkurl{#1}}
\newcommand{\filepath}[1]{\nolinkurl{#1}}
\newcommand{\schemeheader}[3]{
  \begin{center}
    {\LARGE \textbf{#1}}\\[0.5em]
    {\large #2}\\[0.4em]
    {\normalsize \textit{#3}}
  \end{center}
  \vspace{0.8em}
}



\begin{document}
\schemeheader
  {Scheme 2}
  {Stage 2 Abstract-Level Structured Coding Scheme}
  {Prompt-guided abstract coding with deterministic repair and rule-based fallback}

\section*{Purpose}
This scheme standardizes abstract-level extraction for all eligible chronic pain reviews. It is the main corpus-wide coding layer used to describe review characteristics, classify the function of biopsychosocial language, identify the presence of biological, psychological, and social content, flag conceptual problems, and generate a provisional biopsychosocial typology for downstream synthesis.

\section*{Operational Status}
\begin{itemize}
\item Workflow position: Stage 2 abstract coding after Stage 1 eligibility screening.
\item Operational mode: structured LLM-first coding with archived JSON batch outputs, deterministic vocabulary normalization, and rule-based fallback when model output is incomplete or unavailable.
\item Canonical implementation basis: actual Stage 2 output schema in \filepath{src/review_stages/04_extraction/outputs/stage2_abstract_coding.csv}.
\end{itemize}

\section*{Canonical Source Paths}
\begin{itemize}
\item \filepath{src/protocol/codebooks/stage2_codebook.md}
\item \filepath{src/review_stages/04_extraction/codebooks/stage2_codebook.csv}
\item \filepath{src/bps_review/extraction/stage2.py}
\item \filepath{src/bps_review/extraction/llm_stage2.py}
\item \filepath{src/review_stages/04_extraction/outputs/stage2_abstract_coding.csv}
\item \filepath{src/review_stages/04_extraction/outputs/stage2_llm_structured_coding.csv}
\item \filepath{src/review_stages/04_extraction/outputs/llm_stage2_structured_batches.jsonl}
\end{itemize}

\section*{Carried-Through Metadata Fields}
\begin{description}
\item[\field{record_id}] Internal record identifier.
\item[\field{database}] Source database.
\item[\field{pmid} / \field{pmcid} / \field{doi}] External identifiers when available.
\item[\field{title} / \field{abstract}] Primary text used for coding.
\item[\field{year} / \field{journal} / \field{authors} / \field{country_contact_author}] Descriptive metadata retained for synthesis.
\item[\field{publication_types}] PubMed or source publication-type metadata.
\item[\field{objective_text}] Objective sentence or excerpt extracted from the abstract.
\item[\field{screening_status} / \field{screening_reason}] Stage 1 status inherited into the Stage 2 dataset.
\end{description}

\section*{Coded Fields and Controlled Values}
\begin{description}
\item[\field{review_type}] \texttt{systematic review}; \texttt{meta-analysis}; \texttt{network meta-analysis}; \texttt{umbrella review}; \texttt{scoping or mapping review}; \texttt{rapid review}; \texttt{realist review}; \texttt{integrative review}; \texttt{narrative or expert review}; \texttt{other evidence synthesis}; \texttt{unclear}.
\item[\field{objective_category}] \texttt{conceptual}; \texttt{clinical}; \texttt{methodological}; \texttt{epidemiological}; \texttt{mixed}; \texttt{unclear}.
\item[\field{objective_category_source}] Indicates whether the objective classification came from the structured LLM, a repaired LLM batch, or the deterministic fallback.
\item[\field{icd11_pain_category}] \texttt{chronic secondary musculoskeletal pain}; \texttt{chronic neuropathic pain}; \texttt{chronic cancer-related pain}; \texttt{chronic postsurgical or posttraumatic pain}; \texttt{chronic secondary headache or orofacial pain}; \texttt{chronic secondary visceral pain}; \texttt{chronic primary pain}; \texttt{mixed or unspecified chronic pain}; \texttt{unclear}.
\item[\field{musculoskeletal_flag}] \texttt{yes}; \texttt{no}; \texttt{unclear}.
\item[\field{bps_mention_location}] \texttt{title only}; \texttt{abstract only}; \texttt{title and abstract}; \texttt{unclear}.
\item[\field{bps_function}] \texttt{explanatory framework}; \texttt{intervention rationale}; \texttt{organizing principle}; \texttt{justification}; \texttt{background framing}; \texttt{conclusion}; \texttt{policy/practice implication}; \texttt{rhetorical label}; \texttt{unclear}.
\item[\field{bio_mentioned} / \field{psych_mentioned} / \field{social_mentioned}] \texttt{yes} or \texttt{no}.
\item[\field{quality_assessment_reported}] \texttt{yes}; \texttt{no}; \texttt{unclear}.
\item[\field{psychological_concepts_detected}] Normalized list of psychological concepts extracted from title and abstract.
\item[\field{theoretical_frameworks_detected}] Normalized list of named frameworks or model labels.
\item[\field{conceptual_problem_flags}] \field{vague_definition}; \field{tokenistic_bps}; \field{missing_social}; \field{missing_biology}; \field{mechanistic_absence}; \field{construct_overlap}; \field{parallel_listing_without_integration}; \field{none}.
\item[\field{provisional_typology}] \texttt{potential integrative signal}; \texttt{multifactorial signal}; \texttt{pseudo-bps or partial signal}; \texttt{rhetorical label signal}.
\item[\field{stage3_candidate}] \texttt{yes} or \texttt{no}.
\item[\field{stage3_priority}] \texttt{high}; \texttt{medium}; \texttt{low}.
\item[\field{coding_rationale}] One-sentence rationale supporting the coding bundle.
\item[\field{coding_method}] Operational provenance of the row, for example \field{llm_structured}, \field{llm_batch_fallback}, or \field{rule_based_fallback}.
\item[\field{llm_model}] Model identifier when LLM coding was used.
\end{description}

\section*{Prompt Rules Used in the Structured LLM Stage}
\begin{itemize}
\item Code only from title, abstract, publication types, and journal metadata.
\item Do not infer content that is absent from the record.
\item Do not treat the single lexical token \texttt{biopsychosocial} as proof that biological, psychological, and social domains are all substantively present.
\item Use \texttt{explanatory framework} only when the review explicitly treats BPS as a model or framework explaining pain or pain-related disability.
\item Use \texttt{intervention rationale} when BPS language mainly justifies multimodal or interdisciplinary treatment.
\item Use \texttt{organizing principle} when BPS language structures the scope or categories without specifying integration mechanisms.
\item Use \texttt{rhetorical label} when BPS is invoked ceremonially, aspirationally, or symbolically without substantive analytic work.
\item Set \field{stage3_candidate} to \texttt{yes} for musculoskeletal reviews and mixed or unspecified chronic pain reviews where musculoskeletal relevance cannot be ruled out.
\item Use \texttt{potential integrative signal} only when all three core domains are substantively present and the abstract signals cross-domain explanation or organization beyond simple listing.
\item Use \texttt{multifactorial signal} when all three domains are substantively present but mainly parallel.
\item Use \texttt{pseudo-bps or partial signal} when one or more core domains are thin or absent despite BPS language.
\item Use \texttt{rhetorical label signal} when the BPS label is mainly symbolic or confined to framing or conclusion.
\item Return \field{conceptual_problem_flags} as \texttt{none} only when no inferable conceptual problem is present.
\end{itemize}

\section*{Deterministic Repair and Fallback Logic}
\begin{itemize}
\item Missing or malformed model responses are repaired against the fixed vocabularies listed above.
\item Aliases such as \texttt{scoping review} or \texttt{narrative review} are normalized to the implemented labels.
\item \field{stage3_candidate} is forced to \texttt{yes} whenever \field{musculoskeletal_flag} is \texttt{yes} or \texttt{unclear}.
\item Conceptual problem flags are post-derived from typology, BPS function, and missing domains. Rhetorical usage triggers \field{tokenistic_bps} and \field{vague_definition}; parallel non-mechanistic usage triggers \field{parallel_listing_without_integration} and \field{mechanistic_absence}; absent biological or social content triggers \field{missing_biology} or \field{missing_social}.
\item When the LLM stage fails, the rule-based fallback still generates review type, objective category, ICD-11 category, BPS function, domain mention flags, concept detections, and provisional Stage 3 priority.
\end{itemize}

\section*{Implementation Notes and Documented Divergences}
\begin{itemize}
\item The prose codebook lists \texttt{treatment-oriented} as an objective class, but the implemented structured prompt and output schema do not currently use that label.
\item The CSV codebook historically contains \field{spiritual_existential_mentioned} and \field{lifestyle_mentioned}; those fields are not present in the current Stage 2 output file and therefore are not treated as canonical in this dossier.
\item The prose codebook and the executable code both emphasize that BPS language alone is insufficient evidence of true triadic coverage.
\end{itemize}

\section*{Primary Outputs Using This Scheme}
\begin{itemize}
\item \filepath{src/review_stages/04_extraction/outputs/stage2_abstract_coding.csv}
\item \filepath{src/review_stages/04_extraction/forms/stage2_double_code_subset.csv}
\item \filepath{src/review_stages/04_extraction/outputs/stage2_objective_llm_assist.csv}
\item \filepath{src/review_stages/04_extraction/outputs/llm_objective_pilot.json}
\end{itemize}

\end{document}
